% =============================================================================
% FST-Nash: Game-Theoretic Diagnostics for Chaperone Systems
% Version 1.4 — Zenodo Preprint
% Author: Lukas Geiger, Independent Researcher, Bernau
% =============================================================================
\documentclass[12pt,a4paper]{article}

\usepackage[utf8]{inputenc}
\usepackage[T1]{fontenc}
\usepackage[english]{babel}
\usepackage{lmodern}
\usepackage{microtype}
\usepackage{amsmath,amssymb,amsthm}
\usepackage{graphicx}
\usepackage{booktabs}
\usepackage{multirow}
\usepackage{array}
\usepackage{longtable}
\usepackage{tabularx}
\usepackage{needspace}
\usepackage[margin=2.5cm]{geometry}
\usepackage{xcolor}
\usepackage{xurl}
\input{glyphtounicode}
\pdfgentounicode=1
\usepackage[unicode=true,colorlinks=true,linkcolor=blue,citecolor=blue,urlcolor=blue]{hyperref}
\hypersetup{
  pdftitle={Construction-Conditioned Potential-Game Diagnostics for Chaperone Models},
  pdfauthor={Lukas Geiger},
  pdfsubject={Functional Stability Theory; chaperone systems; potential-game diagnostics},
  pdfkeywords={chaperone systems, potential game, Nash equilibrium, symmetry breaking, protein folding, non-equilibrium thermodynamics, game theory, fold switching, functional stability theory},
  pdflang={en-US},
  bookmarksnumbered=true,
  pdfpagemode=UseOutlines
}

\newtheorem{proposition}{Proposition}
\newtheorem{definition}{Definition}
\newcolumntype{L}[1]{>{\raggedright\arraybackslash}p{#1}}
\newcolumntype{Y}{>{\raggedright\arraybackslash}X}
\emergencystretch=3em

\title{Construction-Conditioned Potential-Game\\
Diagnostics for Chaperone Models}

\author{Lukas Geiger\\
\small Independent Researcher, Bernau\\
\small Code: \href{https://github.com/research-line/fst-nash}{research-line/fst-nash}}

\date{October 2026}

\begin{document}
\maketitle

% =============================================================================
\begin{abstract}
We examine $2\times 2$ normal-form games constructed from selected
chaperone--substrate and fold-switching examples with the exact
potential-game (PG) test of Monderer and Shapley (1996). Under the
modeling convention used here, a shared interaction contrast is encoded
as PG and an asymmetric contrast as non-PG. This convention is motivated
by Xu's (2022) symmetry framework, but it is not derived from a
chaperone reaction network and therefore does not establish an
$\mathrm{S4}\iff\mathrm{PG}$ equivalence.

The 16-case table is consequently a \emph{construction ledger}, not an
independently validated biological atlas. XCL1/\allowbreak Lymphotactin
provides an internal salt-shift reconstruction ($23.7\%$~Ltn40 at
$150$\,mM~NaCl) because the same Tyler et~al.\ (2012) data set supplies
the calibration and comparison. The Thermosome example shows that the
same biological system can receive opposite PG labels under different
game constructions. A five-axis validation ledger yields $0/8$ full
validation passes. Nash-equilibrium identities and sensitivity
boundaries are therefore candidate outputs for future preregistered
tests, not validated predictions.

\medskip\noindent
\textbf{Relation to Prior Work.}
This paper supersedes Section~3 (``Game-Theoretic Stability'') of the
overview paper \emph{Functional Stability Theory~III: Biological Stability
and Nash Frustration} (Geiger~2026, DOI: 10.5281/zenodo.20130573).
The Nash-Frustration correlation reported there ($\rho_S = 0.44$,
$p = 0.033$, $n = 24$) is not discriminatory under the earlier
minimum-and-step-size construction: for a positive-definite Hessian,
$\rho(I-\eta H)<1$ follows whenever
$0<\eta<2/\lambda_{\max}(H)$. Those results are therefore not retained
as biological evidence here; the present paper uses the reformulated
framework instead.
\end{abstract}

\clearpage
\begingroup
\small
\setlength{\parskip}{0pt}
\tableofcontents
\endgroup
\clearpage


% =============================================================================
\section{Introduction}
\label{sec:intro}

Molecular chaperones assist protein folding through diverse mechanisms
ranging from passive holdase activity to ATP-driven conformational cycles.
Xu~\cite{Xu2022}, building on Goloubinoff et~al.~\cite{Goloubinoff2018},
classified chaperone mechanisms via four symmetry conditions on the rate
matrix of chaperone--substrate interactions:
\begin{itemize}
  \item[\textbf{S1}] Detailed balance of binding/unbinding
  \item[\textbf{S2}] Substrate-independent dissociation
  \item[\textbf{S3}] Substrate-independent ATPase rate
  \item[\textbf{S4}] Reversibility: every closed chaperone state has a
    reversible path to an open state without consuming chemical energy
\end{itemize}
When all four hold, the system is at equilibrium.
Breaking any condition yields a non-equilibrium
system. We abbreviate these as Goldilocks~(GG) and
non-Goldilocks~(NGG) throughout; this terminology is ours, not Xu's.
We use ``Goloubinoff landscape'' as shorthand for the
equilibrium/non-equilibrium chaperone research program founded by
Goloubinoff et~al.~\cite{Goloubinoff2018} and formalized into the
S1--S4 conditions by Xu~\cite{Xu2022}.
Xu already distinguishes Hsp70 and Hsp90 by the different subsets of
symmetry conditions they break and by their different non-equilibrium
mechanisms. Our narrower question is whether explicitly specified game
representations can provide an auditable re-encoding of such symmetry
patterns.

\subsection{Contribution}

We apply the \emph{potential-game} (PG) property of Monderer and
Shapley~\cite{MondererShapley1996} to explicitly constructed
$2\times 2$ normal-form games. The four-cycle test audits a specified
payoff representation; it does not independently determine Xu's S4
condition. We use the following proposed labels as a descriptive
partition:
\begin{enumerate}
  \item \textbf{Kinetic NGG} (S3 broken, S4 treated as intact; encoded
    with a shared contrast, hence PG by construction): proposed for routes
    in which rates are substrate-dependent but no directed cycle is
    assumed to shift the equilibrium.
  \item \textbf{Thermodynamic NGG} (S3 and S4 treated as broken; encoded
    with asymmetric contrasts, hence non-PG by construction): proposed for
    routes in which a directed thermodynamic cycle is assumed to shift the
    equilibrium; the present models do not compute cycle affinities.
\end{enumerate}

We document 16~model constructions, including four post-hoc
convention-consistency cases, reconstruct the calibrated XCL1 salt
response, and illustrate construction-conditioning via the Thermosome
meta-demonstration.

\subsection{Relation to Prior Work}

This paper replaces the Nash-Frustration analysis in FST-III~\cite{GeigerFSTIII2025}. That analysis modeled each residue as a player
choosing backbone dihedral angles $(\varphi_i, \psi_i)$ on the
torus~$T^2$, with pairwise cosine couplings. The spectral radius
$\rho(J) < 1$ of the Jacobian $J = I - \eta H$ was interpreted as
``Nash stability.'' However, a random-parameter control
(Section~\ref{sec:tautology}) showed that for a positive-definite
Hessian, $\rho(I-\eta H)<1$ follows under the step-size condition
$0<\eta<2/\lambda_{\max}(H)$. The test is therefore not discriminatory
under that minimum-and-step-size construction. Prior results, including
the Spearman correlation $\rho_S = 0.44$ against B-factors, are not
retained as biological evidence here.

The present framework sets aside the continuous dihedral-angle approach
entirely and instead formulates chaperone--substrate interactions as
discrete $2\times 2$ normal-form games, where the PG test is a well-defined audit of the specified payoff representation.


% =============================================================================
\section{Theory: Four-Cycle Test and the S4-Motivated Construction}
\label{sec:theory}

\begin{definition}[Potential Game, Monderer and Shapley 1996]
A finite $n$-player game is a \emph{potential game} if there exists a
function $\Phi: S_1 \times \cdots \times S_n \to \mathbb{R}$ such that
for every player~$k$, every strategy profile $s_{-k}$, and every pair
of strategies $s_k, s_k'$:
\begin{equation}
  u_k(s_k', s_{-k}) - u_k(s_k, s_{-k})
  = \Phi(s_k', s_{-k}) - \Phi(s_k, s_{-k}).
  \label{eq:pg_def}
\end{equation}
\end{definition}

For a $2\times 2$ game with players $H$ (chaperone) and $S$ (substrate),
strategies $\{0, 1\}$, and payoff matrices $u_H$, $u_S$, the PG
condition reduces to a single scalar test:

\begin{proposition}[Four-Cycle Test for $2\times 2$ Games]
\label{prop:four_cycle}
Define the \emph{interaction index} for each player~$k$:
\begin{equation}
  I_k = u_k[0,0] - u_k[0,1] - u_k[1,0] + u_k[1,1].
  \label{eq:interaction}
\end{equation}
The game is a potential game if and only if $I_H = I_S$.
\end{proposition}

\begin{proof}
Consider the two unilateral paths from $(0,0)$ to $(1,1)$. Along the
first path, $H$ changes the row and $S$ then changes the column; along
the second, $S$ changes the column and $H$ then changes the row. A
scalar exact potential is path independent if and only if
\begin{align}
 &[u_H(1,0)-u_H(0,0)]+[u_S(1,1)-u_S(1,0)] \nonumber\\
 ={}&[u_S(0,1)-u_S(0,0)]+[u_H(1,1)-u_H(0,1)].
\end{align}
Rearranging gives $I_H=I_S$, the four-cycle condition for this
orientation. The converse follows by constructing $\Phi$ along either
path~\cite{MondererShapley1996}.
\end{proof}

\subsection{Construction Rule Motivated by Xu's S4}

Xu's fourth symmetry condition~(S4)~\cite{Xu2022} requires that every
closed chaperone state has a reversible path to an open state without
consuming chemical energy---i.e., that the chaperone--substrate
reaction network is microscopically reversible.

In the present models, a case treated as S4-intact is represented with a
shared interaction contrast, which makes $I_H=I_S$ by construction. A
case treated as S4-broken is represented with asymmetric contrasts,
which generally makes $I_H\neq I_S$. These are modeling rules, not a
derivation from transition rates or a proof of biological equivalence.

\textbf{Construction dependence.} Potentiality is a property of the
specified game. Player partition, coarse-graining, utility definition,
and cardinal scale can change the PG result. A coarse-grained four-state
game can also hide dissipative cycles. Thus PG does not imply biological
reversibility or S4 unless a complete network-to-game mapping and its
cycle affinities are independently established. The Thermosome
meta-demonstration (Section~\ref{sec:thermosome}) illustrates this
dependence.


% =============================================================================
\section{Methods}
\label{sec:methods}

\subsection{Game Construction}
\label{sec:game_construction}

Each chaperone--substrate system is modeled as a $2\times 2$ normal-form
game:
\begin{itemize}
  \item \textbf{Player~$H$} (chaperone): strategies
    $\{\text{Inactive/Resting}, \text{Active/Engaged}\}$
  \item \textbf{Player~$S$} (substrate): strategies
    $\{\text{Unfolded/Free}, \text{Folded/Bound}\}$
\end{itemize}

The payoff matrix for player~$k \in \{H, S\}$ is:
\begin{equation}
  u_k = \begin{pmatrix}
    u_k[0,0] & u_k[0,1] \\
    u_k[1,0] & u_k[1,1]
  \end{pmatrix},
\end{equation}
where larger entries are treated computationally as preferred model
scores. Some terms are energy-derived and reported in kcal/mol; other
terms are transformations of kinetic ratios or assumptions. The current
case models do not provide a universal thermodynamic mapping from
$K_D$, $k_\text{cat}$, and conformational energies onto one common
cardinal utility scale. Accordingly, exact-PG differences in kcal/mol
are construction-level score contrasts, not direct physical
measurements. A claim-bearing thermodynamic model would instead require
matched forward/reverse rates or explicitly defined utilities such as
$u=-G$ at fixed temperature and standard state.

The coupling energy $\Delta_k$ represents the change in player~$k$'s
payoff when the other player switches strategy:
\begin{equation}
  \Delta_k = u_k[\text{coupled}] - u_k[\text{uncoupled}].
\end{equation}

\subsection{Four-Cycle Test (PG Test)}

For each calibrated system, we compute $I_H$ and $I_S$ via
Eq.~\eqref{eq:interaction} and evaluate:
\begin{equation}
  \text{PG violation} = |I_H - I_S|.
\end{equation}
If the violation is below a numerical threshold
($\varepsilon = 10^{-10}$), the specified game is classified as an
exact PG. This numerical identity is not an empirical S4 decision and
does not propagate measurement uncertainty.

\subsection{Nash Equilibrium Search}

Pure Nash equilibria are identified by checking all four strategy
profiles for unilateral deviation incentives. A profile $(s_H^*, s_S^*)$
is a pure NE if:
\begin{align}
  u_H(s_H^*, s_S^*) &\geq u_H(s_H', s_S^*) \quad \forall\, s_H', \\
  u_S(s_H^*, s_S^*) &\geq u_S(s_H^*, s_S') \quad \forall\, s_S'.
\end{align}

\subsection{Regime Classification}

Based on PG test and Xu's symmetry conditions~\cite{Xu2022}:
\begin{enumerate}
  \item \textbf{GG} (Equilibrium): S1--S4 treated as intact; the chosen
    shared-contrast representation is PG; no ATP or no modeled
    substrate stimulation
  \item \textbf{Kinetic NGG}: S3 broken, S4 treated as intact; the
    chosen shared-contrast representation is PG
  \item \textbf{Thermodynamic NGG}: S3 and S4 treated as broken; the
    chosen asymmetric-contrast representation is non-PG
  \item \textbf{Special cases}: non-standard $2\times 2$ structure
    (3-player, oscillatory, regulated)
\end{enumerate}

\subsection{Symmetry Mapping}

For each system, we evaluate S1--S4 based on published experimental
data:
\begin{itemize}
  \item S1: Does the system obey detailed balance for binding/unbinding?
  \item S2: Is dissociation substrate-independent?
  \item S3: Is the ATPase rate substrate-independent?
  \item S4: Does the cited biological route meet Xu's reversibility
    condition? This label is recorded separately from the game result;
    the present construction encodes an S4-intact label with a shared
    contrast and therefore with $I_H=I_S$.
\end{itemize}


% =============================================================================
\section{Results}

% --- Result 1 ----------------------------------------------------------------
\subsection{Construction-Conditioned Regime Proposal}
\label{sec:regime_trinity}

The 16~game constructions can be organized into three proposed labels
(Table~\ref{tab:atlas}). This is a descriptive grouping of the chosen
representations, not an independently learned biological partition:

\begin{enumerate}
  \item \textbf{Equilibrium (GG):} 6~systems. PG $=$ True,
    no S3-breaking. Includes XCL1 (metamorphic fold-switching),
    Prefoldin (holdase), and 4~post-hoc convention-consistency systems
    (DnaJ, Hsp26, Hsp33, Hsc70).

  \item \textbf{Kinetic NGG:} 2~systems (Hsp70/DnaJ, SecA).
    Xu identifies Hsp70-mediated folding as breaking S3. His necessary
    condition for breaking S4 requires at least two ATP molecules to be
    hydrolyzed sequentially per cycle, which is absent from the Hsp70
    route considered here. The biological S4 label is therefore
    consistent with Xu. Separately, PG$=$True in our matrix follows
    from the assumed shared interaction contrast; it is not an
    independent test of that label.

  \item \textbf{Thermodynamic NGG:} 4~systems (Hsp90/Cdc37, GroEL/GroES,
    ClpB/Hsp104, p97/VCP). PG $=$ False with constructed score
    differences of $0.5$--$2.5$. These systems are represented as breaking both S3 and S4. The motivating
    hypothesis is that co-chaperones, multi-ring architectures, or domain
    asymmetries create directed thermodynamic cycles that shift the
    equilibrium; the present $2\times 2$ scores encode this assumption and
    do not compute or test cycle affinities. For GroEL, substrate release requires sequential ATP
    hydrolysis in the trans-ring, violating Xu's S4~\cite{Xu2022}.
\end{enumerate}

Four additional systems (Trigger Factor, KaiB/KaiC, Hsp70$\to$Hsp90
handoff, MAD2) fall outside the standard three-regime classification due
to non-standard game structures (ribosome-coupled, oscillatory,
3-player, or regulated metamorphic).


% --- Result 2 ----------------------------------------------------------------
\clearpage
\subsection{Construction Ledger}
\label{sec:atlas}

\begingroup
\footnotesize
\setlength{\tabcolsep}{3pt}
\renewcommand{\arraystretch}{1.12}
\begin{longtable}{@{}rL{0.23\textwidth}cccL{0.13\textwidth}L{0.33\textwidth}@{}}
\caption{Construction Ledger: 16 specified game representations audited
with the exact-PG test. The final column records input status and the
construction that determines PG or non-PG. No row is an independent
validation of a biological label (see Section~\ref{sec:thermosome}).
NE $=$ number of pure Nash equilibria in the specified game.}
\label{tab:atlas}
\\
\toprule
\# & System & ATP? & PG & NE & Regime & Evidence / PG basis \\
\midrule
\endfirsthead
\toprule
\# & System & ATP? & PG & NE & Regime & Evidence / PG basis \\
\midrule
\endhead
\bottomrule
\endfoot
\bottomrule
\endlastfoot
\multicolumn{7}{@{}l}{\textit{Regime 1 --- Equilibrium (GG)}} \\
1  & XCL1 (metamorphic) & No & $\checkmark$ & \textbf{2} & GG
   & Symmetric game by construction; calibrated. \\
2  & Prefoldin (holdase) & No & $\checkmark$* & \textbf{2}* & GG*
   & Estimated parameters ($f{=}0.5$); not independently calibrated. \\
3  & DnaJ alone & No & $\checkmark$ & 1 & GG
   & Shared-$\Delta$ convention; post-hoc consistency case. \\
4  & Hsp26 sHsp & No & $\checkmark$ & 1 & GG
   & Shared-$\Delta$ convention; post-hoc consistency case. \\
5  & Hsp33 redox & No & $\checkmark$ & 1 & GG
   & Shared-$\Delta$ convention; post-hoc consistency case. \\
6  & Hsc70 alone & Basal & $\checkmark$ & 1 & GG
   & Shared-$\Delta$ convention; post-hoc consistency case. \\
\midrule
\multicolumn{7}{@{}l}{\textit{Regime 2 --- Kinetic NGG}} \\
7  & Hsp70/DnaJ & Yes & $\checkmark$ & 1 & kin.\ NGG
   & Shared-$\Delta$ from experimental $K_D$ values; calibrated. \\
8  & SecA & Yes & $\checkmark$ & 1 & kin.\ NGG
   & Shared-$\Delta$ from experimental $K_D$ values; calibrated. \\
\midrule
\multicolumn{7}{@{}l}{\textit{Regime 3 --- Thermodynamic NGG}} \\
9  & Hsp90/Cdc37 & Yes & $\times$ & 1 & therm.\ NGG
   & Asymmetric $\Delta$ (calibrated/estimated); calibrated. \\
10 & GroEL/GroES & Yes & $\times$ & 1 & therm.\ NGG
   & Experimentally motivated ring asymmetry; calibrated. \\
11 & ClpB/Hsp104 & Yes & $\times$ & 1 & therm.\ NGG
   & Threading asymmetry; calibrated. \\
12 & p97/VCP & Yes & $\times$ & 1 & therm.\ NGG
   & Estimated D1/D2 asymmetry; calibrated. \\
\midrule
\multicolumn{7}{@{}l}{\textit{Special cases (non-standard regime)}} \\
13 & Trigger Factor & No & $\times$ & 1 & non-GG
   & Ribosome asymmetry using experimental $K_D$; calibrated. \\
14 & KaiB/KaiC & Indirect & $\times$ & 1 & non-GG
   & Oscillatory system with experimental kinetics; calibrated. \\
15 & Hsp70$\to$Hsp90 (3-player) & Yes & $\times$ & 1 & non-PG
   & Three-player model, not a $2\times 2$ game; calibrated. \\
16 & MAD2 (regulated metamorphic) & No$^\ddagger$ & $\times$ & 1 & non-PG
   & TRIP13-regulated asymmetry; calibrated. \\
\end{longtable}

\noindent{\footnotesize
\textbf{Honesty note.} All PG$=$True systems have PG as an algebraic
consequence of the modeling convention: XCL1 from a symmetric game
(identical players), Hsp70/DnaJ and SecA from the assumption
$\Delta_H = \Delta_S$ (experimental $K_D$ values, but symmetric coupling
is \emph{assumed}), and post-hoc cases from the shared-$\Delta$
convention.
All non-PG systems have non-PG from $\Delta_H \neq \Delta_S$ or from
asymmetric game structure. PG/non-PG is thus a \emph{convention decision
during modeling}, not an emergent empirical result. * Prefoldin is based
on estimated parameters ($K_D$, $\Delta G_\text{conf}$ unpublished).
$^\ddagger$ MAD2 has no direct ATP; TRIP13 AAA+ ATPase resets MAD2.}
\endgroup


% --- Result 3 ----------------------------------------------------------------
\subsection{XCL1 Salt-Shift: Internal Calibrated Reconstruction}
\label{sec:xcl1}

XCL1/Lymphotactin is a metamorphic protein that interconverts between a
chemokine fold (Ltn10) and a $\beta$-sheet fold (Ltn40) with
$K_\text{eq} = 1.2$ at $25\,{}^\circ\text{C}$, $0$\,mM~NaCl as the
project's calibration anchor. The internal records refer to both
Tuinstra et~al.~\cite{Tuinstra2008} and Tyler et~al.~\cite{Tyler2012};
the precise primary-text attribution of this value remains unverified. We model XCL1 as a symmetric $2\times 2$
coordination game between two identical monomers, with strategies
$\{\text{Chemokine}, \text{Beta-sheet}\}$.

The game yields 2~pure NE: (Chem, Chem) $=$ Ltn10 and
(Beta, Beta) $=$ Ltn40, representing the metamorphic bistability as a
coordination equilibrium.

\subsubsection{Salt-Dependent Population Shift}

Tyler et~al.~\cite{Tyler2012} showed that NaCl shifts the population
from ${\sim}55\%$ Ltn40 ($0$\,mM) to ${\sim}24\%$ Ltn40 ($150$\,mM) via
chloride binding at Arg23/Arg43 ($K_d \approx 16$\,mM from CSP
titration). We model this as a chloride-binding saturation:

\begin{equation}
  f(\text{Cl}^-) = \frac{[\text{Cl}^-]}{[\text{Cl}^-] + K_d}, \qquad
  \Delta\Delta G([\text{NaCl}])
  = \Delta\Delta G_{150}
    \frac{f([\text{NaCl}])}{f(150\,\text{mM})},
  \label{eq:salt_model}
\end{equation}
with $\Delta\Delta G_{150} = 0.8$\,kcal/mol (Tyler 2012 ZZ-exchange:
$\Delta\Delta G(k_f)=0.5$ and $\Delta\Delta G(k_r)=0.3$\,kcal/mol,
summing to $0.8$).
Thus $0.8$\,kcal/mol is the encoded 150-mM anchor, not the asymptotic
shift; the implied asymptote is approximately $0.885$\,kcal/mol. This
normalization matches the implementation and table. Its precise source
semantics and uncertainty still require a primary-full-text audit.

\begin{table}[!htbp]
\centering
\caption{XCL1 Salt-Shift Reconstruction. Model calibrated at
$K_\text{eq}(0\,\text{mM}) = 1.2$, $\Delta\Delta G_{150} = 0.8$\,kcal/mol,
and $K_d = 16$\,mM.}
\label{tab:xcl1_salt}
\begin{tabular}{@{}rcccc@{}}
\toprule
{[NaCl]} (mM) & $f(\text{Cl}^-)$ & $K_\text{eq}$ & \% Ltn40 & Experimental \\
\midrule
0   & 0.00 & 1.200 & 54.5 & 54.5\% (project anchor) \\
10  & 0.39 & 0.676 & 40.3 & --- \\
50  & 0.76 & 0.387 & 27.9 & --- \\
75  & 0.82 & 0.351 & 26.0 & Tyler: ``similar to 150\,mM''  \\
150 & 0.90 & 0.311 & 23.7 & ${\sim}24\%$ Tyler ZZ-exch.\  \\
300 & 0.95 & 0.291 & 22.5 & --- \\
365 & 0.96 & 0.287 & 22.3 & CSP saturation (Tyler) \\
\bottomrule
\end{tabular}
\end{table}

\subsubsection{Independence Status}

The $150$\,mM reconstruction ($23.7\%$) uses Tyler's
$\Delta\Delta G$ as input and recalculates the same evidence family.
It is not an independent prediction. The remaining internal candidate
outputs are:

\begin{enumerate}
  \item \textbf{75\,mM grid point:} The model gives $26.0\%$. Tyler's
    qualitative statement that $70$\,mM and $150$\,mM had a similar
    effect comes from the calibration evidence family and is not an
    external test.
  \item \textbf{Saturation curve shape:} The nonlinear form is
    determined by chloride-binding saturation ($K_d = 16$\,mM from CSP),
    not derivable from Tyler's population data alone.
  \item \textbf{K25A mutant:} The model output is
    $K_\text{eq} = 0.40$ ($28.6\%$~Ltn40). It is not treated here as an
    independently preregistered mutant prediction.
\end{enumerate}


% --- Result 4 ----------------------------------------------------------------
\subsection[Construction-Conditioning: Thermosome]{Construction-Conditioning: Thermosome\\ Meta-Demonstration}
\label{sec:thermosome}

To demonstrate that PG/non-PG depends on the game formulation and is
not fixed by the biological system alone, we analyze the Thermosome of
\emph{Thermoplasma acidophilum}---a Group~II chaperonin with
$(\alpha\beta)_4(\alpha\beta)_4$ hexadecameric architecture and negative
inter-ring cooperativity~\cite{Bigotti2008,Gutsche2000}---using two deliberately
contrasting game encodings:

\begin{table}[!htbp]
\centering
\caption{Thermosome: Same Biological System, Two Game Models, Opposite
PG Results.}
\label{tab:thermosome}
\footnotesize
\begin{tabularx}{\textwidth}{@{}L{0.22\textwidth}YccrrL{0.13\textwidth}@{}}
\toprule
Model & Players & PG & NE & $I_1$ & $I_2$ & Regime \\
\midrule
A: Ring vs.\ Ring & $\alpha_4\beta_4$-Ring 1 vs.\ Ring 2 &
  $\checkmark$ & 2 & $-12.92$ & $-12.92$ & GG \\
B: Chaperone vs.\ Substrate & Hexadecamer vs.\ Substrate &
  $\times$ & 1 & $-4.0$ & $-1.5$ & therm.\ NGG \\
\bottomrule
\end{tabularx}
\end{table}

\textbf{Model~A} treats the two rings as identical players in a symmetric
anti-coordination game $\{$Tight, Relaxed$\}$. PG follows algebraically
because $u_1 = u_2^\top$ (identical players $\Rightarrow I_1 = I_2$ by
construction). Its two static NE are (Tight, Relaxed) and
(Relaxed, Tight). They represent complementary ring states but do not
by themselves establish temporal alternation; that would require a
transition or Markov model.

\textbf{Model~B} treats the hexadecamer and substrate as asymmetric
players. It converts an ATP-binding input and a one-way substrate-folding
acceleration into score terms $\Delta_H=4.0$ and $\Delta_S=1.5$.
Because those observables do not by themselves define a shared
thermodynamic scale, their $2.5$ difference is a constructed PG-score
contrast rather than a measured free-energy violation.

The pair illustrates model dependence, not two equally validated
biological mechanisms. The choice is a modeling decision and motivates
construction-conditioning transparency throughout this paper.

\subsubsection{Robustness}

A scan over $\Delta_H \in [0.5, 8.0] \times \Delta_S \in [0.5, 8.0]$
confirms that Model~B yields PG$=$True exclusively on the diagonal
$\Delta_H = \Delta_S$, as follows algebraically from the payoff
structure. This is sensitivity to the specified score scan; the scan
does not establish which parameter combinations are physically
plausible.


% =============================================================================
\section{Known Limitations and Tautology Transparency}
\label{sec:tautology}

\subsection{What Is Tautological}

\begin{enumerate}
  \item \textbf{PG/non-PG follows from payoff construction:}
    $\Delta_H = \Delta_S \Rightarrow$ PG (algebraically);
    $\Delta_H \neq \Delta_S \Rightarrow$ non-PG (algebraically).
    The PG test tests the payoff symmetry, not a system property.

  \item \textbf{NE multiplicity:} In a generic (tie-free) $2\times 2$
    exact potential game, two pure NE occur if and only if neither player
    has a strictly dominant strategy; this is a consequence of the game
    structure, not a biological finding.

  \item \textbf{Prior Nash-Frustration results (FST-III v3):} The
    spectral radius $\rho(I-\eta H)<1$ at a non-degenerate minimum with
    positive-definite Hessian follows for
    $0<\eta<2/\lambda_{\max}(H)$. Under that construction it is not a
    discriminator of biological stability. Random parameters satisfy the same criterion
    ($\rho(J)<1$ in $50/50$ random and $10/10$ fitted runs).
\end{enumerate}

\noindent\textbf{Scale limitation.} Exact potentiality is a cardinal
property. Separate positive rescalings of player utilities can preserve
best responses and Nash equilibria while changing $I_H=I_S$. A common
physical scale must therefore be derived before exact PG is interpreted
thermodynamically; weighted and ordinal potentials answer weaker,
different questions.

\subsection{Candidate Outputs Not Yet Independently Validated}

\begin{enumerate}
  \item \textbf{NE identity and sensitivity:} The specified matrices
    produce equilibrium profiles and score-boundaries that can be frozen
    for future tests.
  \item \textbf{Salt response:} The calibrated XCL1 reconstruction
    produces a concentration-response curve, but no untouched
    concentration measurement has yet been tested after model freeze.
  \item \textbf{Regime labels:} The proposed grouping is an auditable
    re-encoding of selected symmetry patterns, not an empirically
    discovered ontology.
\end{enumerate}

\subsection{Pre-Review Validation-Evidence Gate}

The machine-readable validation ledger generated by
\path{scripts/validation_evidence_ledger.py} separates eight
headline case routes across five axes: construction conditioning,
experimental dynamics, an AlphaFold baseline, an energetic-frustration
baseline, and an independent chaperone-cycle anchor. All eight rows are
construction-conditioned; none contains a computed AlphaFold baseline,
an energetic-frustration baseline, or a quantitative independent cycle
anchor. Consequently, the full-stack validation count is $0/8$.

This gate does not erase the narrower XCL1 population consistency check,
but it preserves its partial-circularity caveat. GroEL remains only a
qualitative partial anchor: two of three kinetic categories are
consistent, while all modeled cycle times miss the stated
$15$--$20$\,s experimental range (primary source pending verification). The current evidence therefore supports
diagnostic classification only, not a validated mechanistic model or a
claim upgrade.


% =============================================================================
\section{Discussion}
\label{sec:discussion}

\subsection{Honest Balance}

We evaluate our framework against 13~open questions in the chaperone
field (Table~\ref{tab:balance}). The framework addresses ${\sim}38\%$ of
questions (5/13), of which only one (Q9, XCL1 salt-shift) has a
quantitative, partially circular consistency check.

\begin{table}[!htbp]
\centering
\caption{Honest Balance: Framework Coverage of 13~Open Questions.}
\label{tab:balance}
\begin{tabular}{@{}llc@{}}
\toprule
Category & Questions & Fraction \\
\midrule
QUANTITATIVE CHECK (partly circular) & Q9 (XCL1 salt-shift) & 7.7\% \\
MODEL OUTPUT (construction-conditioned) & Q1, Q6, Q10, Q13 & 30.8\% \\
RE-DESCRIBE & Q4, Q8 & 15.4\% \\
NONE & Q2, Q3, Q5, Q7, Q11, Q12 & 46.2\% \\
\bottomrule
\end{tabular}
\end{table}

${\sim}46\%$ of questions lie entirely outside the framework's scope.
This balance is the honest foundation of this paper.

\subsection{Comparison with Existing Approaches}

The PG test provides a \emph{diagnostic} classification layer that
complements rather than replaces existing approaches:

\begin{itemize}
  \item \textbf{Statistical mechanics (Ising, MWC):} If the exact potential
    of a PG game is identified with a negative free energy on a common
    physical scale, pure-NE search reduces to locating free-energy minima
    under unilateral moves and adds nothing beyond statistical mechanics.
    Within a specified cardinal representation, PG is necessary but not
    sufficient for this reduction; the present scores do not establish
    such a scale (Section~\ref{sec:game_construction}).

  \item \textbf{AlphaFold / structural prediction:} Our framework
    operates on a different level---it labels a specified game representation
    of conformational regimes, not the structure itself. For metamorphic
    proteins (XCL1, MAD2), the NE analysis supplies
    construction-conditioned candidate descriptions of bistability.

  \item \textbf{Xu/Goloubinoff framework:} Xu already distinguishes
    mechanisms by their broken symmetry subsets. Our contribution is a
    proposed game-theoretic re-encoding and audit, not the discovery of
    that biological distinction.
\end{itemize}

\subsection{What This Paper Does Not Claim}

\begin{itemize}
  \item Nash equilibrium ``explains'' chaperone function (it classifies).
  \item PG/non-PG is an ``explanation'' for a phenomenon (it is
    construction-conditioned).
  \item ``Three regimes instead of two'' as an extension of Goloubinoff
    (it is substructure).
  \item AF2 failure is ``solved'' by Nash analysis (it is reformulated).
  \item ATP costs are ``predicted'' by Nash analysis (Goloubinoff
    already explains them).
  \item The $150$\,mM population is an independent prediction (it is a
    calibrated internal reconstruction).
\end{itemize}


% =============================================================================
\section{Conclusion}
\label{sec:conclusion}

The four-cycle test is mathematically valid for a specified game, but
its biological interpretation is conditional on player partition,
coarse-graining, utility definition, and cardinal scale. Current
evidence does not establish an $\mathrm{S4}\iff\mathrm{PG}$ equivalence, a
validated 16-system atlas, or an independent XCL1 prediction. The
present contribution is a construction ledger with explicit failure
modes and a $0/8$ validation status.

The next claim-bearing step is a preregistered network-to-game
construction from matched forward and reverse rates, with uncertainty
propagation, a hidden-cycle audit, and a blinded external endpoint.

All scripts (36~Python files) and result files are available at
\href{https://github.com/research-line/fst-nash}{GitHub: research-line/fst-nash}.


% =============================================================================
\section*{Data and Code Availability}

All computational scripts are located in \texttt{scripts/} (36~files)
with results in \texttt{results/} (JSON format). The template functions
for the four-cycle test, Nash-equilibrium search, and symmetry mapping
are in \texttt{scripts/fold\_switch\_template.py}. Hold-out scripts are
in \texttt{scripts/extension\_B/}. The pre-review evidence gate is
generated by \path{scripts/validation_evidence_ledger.py}.


% =============================================================================
\begin{sloppypar}
\section*{AI Disclosure}

\noindent
AI systems were used extensively in all phases of this work: conception
and research discussion, literature and source work, mathematical analysis
and proof development, creation and checking of the computational scripts,
execution and evaluation of the computational runs, text, translation, and
typesetting.

\medskip

\noindent
Quality assurance: use of several independent models from different
providers with counter-checks and adversarial AI reviews; traceable,
versioned computational scripts with stored result files and checksums;
continuously maintained proof and register documents; strict claim hygiene
(assertions only with documented evidence, addenda instead of silent
corrections); changelogs across all versions.

\end{sloppypar}


% =============================================================================
\begin{thebibliography}{99}

\bibitem{Xu2022}
H.~Xu,
``Non-equilibrium protein folding and activation by ATP-driven
chaperones,''
\textit{Biomolecules}\ \textbf{12}, 832 (2022), DOI: 10.3390/biom12060832.

\bibitem{Goloubinoff2018}
P.~Goloubinoff, A.~S.~Sassi, B.~Fauvet, A.~Barducci, and P.~De~Los~Rios,
``Chaperones convert the energy from ATP into the nonequilibrium
stabilization of native proteins,''
\textit{Nat.\ Chem.\ Biol.}\ \textbf{14}, 388--395 (2018), DOI: 10.1038/s41589-018-0013-8.

\bibitem{MondererShapley1996}
D.~Monderer and L.~S.~Shapley,
``Potential Games,''
\textit{Games Econ.\ Behav.}\ \textbf{14}, 124--143 (1996), DOI: 10.1006/game.1996.0044.

\bibitem{GeigerFSTIII2025}
L.~Geiger,
``Functional Stability Theory~III: Biological Stability and Nash
Frustration,''
Zenodo preprint (2026), DOI: 10.5281/zenodo.20130573.

\bibitem{Tuinstra2008}
R.~L.~Tuinstra, F.~C.~Peterson, S.~Kutlesa, E.~S.~Elgin,
M.~A.~Kron, and B.~F.~Volkman,
``Interconversion between two unrelated protein folds in the
lymphotactin native state,''
\textit{Proc.\ Natl.\ Acad.\ Sci.\ USA}\ \textbf{105}, 5057--5062 (2008), DOI: 10.1073/pnas.0709518105.

\bibitem{Tyler2012}
R.~C.~Tyler, J.~C.~Wieting, F.~C.~Peterson, and B.~F.~Volkman,
``Electrostatic optimization of the conformational energy landscape in
a metamorphic protein,''
\textit{Biochemistry}\ \textbf{51}, 9067--9075 (2012), DOI: 10.1021/bi300842j.

\bibitem{Bigotti2008}
M.~G.~Bigotti and A.~R.~Clarke,
``Chaperonins: The hunt for the Group~II mechanism,''
\textit{Arch.\ Biochem.\ Biophys.}\ \textbf{474}, 331--339 (2008), DOI: 10.1016/j.abb.2008.03.015.

\bibitem{Gutsche2000}
I.~Gutsche, O.~Mihalache, and W.~Baumeister,
``ATPase cycle of an archaeal chaperonin,''
\textit{J.\ Mol.\ Biol.}\ \textbf{300}, 187--196 (2000), DOI: 10.1006/jmbi.2000.3833.

\end{thebibliography}

\end{document}
